\textbf{COLUNI 2021 -- Questão 09}\par\medskip O uso das diferentes grafias dos \underline{porquês} fundamenta-se nas relações sintático-semânticas e se constituem, na escrita, como marca de efeito de sentido.

\par

Com base nessa informação, assinale a alternativa CORRETA, em que está presente a relação de Causa e Consequência:\begin{enumerate}
\item “Eu, dentro de mim, achava que a vovó deveria abrir a porta e explicar ao monge o que acontecera. Mas ela
       sentiu vergonha, não sei \textbf{por quê}.” (LAURITO, Ilka Brunhilde. \textbf{A Menina que fez a América}. São Paulo, FTD, 2019. p. 34)
\item “Um dia, vovó Catarina me disse que possivelmente iria demorar, \textbf{porque} ia à igreja, não à de Sam
       Leone, o padroeiro da cidade, mas à de Sant’Anna, na cidade baixa.” (LAURITO, Ilka Brunhilde. A Menina que fez
       a América. São Paulo, FTD, 2019. p. 73)
\item “Só \textbf{porque}, fazendo cócegas no seu nariz, ele tinha uma vassoura de pelos que eu, por ter nascido
       menina, nunca teria?” (LAURITO, Ilka Brunhilde. \textbf{A Menina que fez a América}. São Paulo, FTD, 2019. pp. 55-56)
\item “Ia fazer promessa. Não lhe perguntei qual era, \textbf{porque} promessa é segredo que só se conta a Deus e
       aos santos.” (LAURITO, Ilka Brunhilde. \textbf{A Menina que fez a América}. São Paulo, FTD, 2019. p.73)\end{enumerate}
